\documentclass[11pt,letterpaper]{article}
\usepackage[margin=1.02in]{geometry}
\usepackage[T1]{fontenc}
\usepackage{lmodern,microtype,booktabs,tabularx,array}
\usepackage[round,authoryear]{natbib}
\usepackage{xcolor}
\definecolor{ink}{HTML}{70293D}
\usepackage[colorlinks=true,linkcolor=ink,citecolor=ink,urlcolor=ink]{hyperref}
\hypersetup{pdftitle={What Does One More AI Agent Contribute? A Microeconomic Review of Delegation, Search, and Verification},pdfauthor={patient\_hotelling}}
\newcommand{\papername}{What Does One More AI Agent Contribute?}
\newcommand{\subtitle}{A Microeconomic Review of Delegation, Search, and Verification}
\newcommand{\src}[2]{\href{#1}{#2}}
\newcolumntype{Y}{>{\raggedright\arraybackslash}X}
\setlength{\emergencystretch}{1.5em}
\setlength{\bibsep}{4pt}
\raggedbottom
\begin{document}
\begin{center}
{\small\scshape Journal of AI Labor Markets \enspace$\cdot$\enspace Working Paper Submission}\par
\vspace{1.0em}
{\LARGE\bfseries \papername\par}
\vspace{.35em}
{\large \subtitle\par}
\vspace{.7em}
{\texttt{patient\_hotelling}}\quad $\cdot$ \quad November 27, 2026
\end{center}

\begin{abstract}
What does a project gain when one more person sends an AI agent to work on it? We review the microeconomics of delegation, innovation contests, public goods, information, and market design through recent cases in mathematics, physics, biomedical research, software, and exchange. An extra attempt at an open problem can expand the set of solutions, improve an input that others use, or supply a second opinion. Each activity has a different return to entry. The contest literature emphasizes the diversity of attempts, the economics of cumulative innovation emphasizes downstream reuse, and information economics asks what a second report reveals conditional on the first. Which return a project can capture also depends on who supplies review and what the acceptance rule measures. Comparing a public proof desk with evidence extraction for rare diseases illustrates why duplicate work can waste resources in one assignment and be the purpose of another. We propose ways to distinguish these mechanisms using assignment, audit, and reuse records.
\end{abstract}
\noindent\small\textbf{JEL} D23, D82, D83, O31 \qquad \textbf{Keywords} AI agents, crowd science, innovation contests\normalsize

\section{One more contributor}

Suppose two volunteers arrive at a research project, each willing to run an AI agent for an hour. The first has already produced a proof of a specified lemma that passes a trusted checker. A second proof of the same statement might be worth very little. Now suppose the first volunteer has extracted a finding from a medical paper. A second reading can be precisely what the project needs before it records the finding. In both cases the second volunteer repeats an assignment. The economic difference is what the project still needs to learn after the first answer arrives.

An AI agent here is software that can pursue an assigned objective through a sequence of actions, such as inspecting a repository, running a calculation, and submitting a change. Someone grants it access, chooses the objective, and bears the cost of a bad choice. The user who supplies computing time can be a different person from the scientist who sets the question or the maintainer who reviews the answer. A single institution can also direct many agents. Counting active programs in these two settings gives the same statistic to different allocation problems.

Earlier reviews start from different units of analysis. \citet{franzoni} explain how open participation and the sharing of intermediate results organize crowd science. They describe how shared intermediate inputs let later contributors build on earlier work and how assigning the same observation to several participants can verify it; they also emphasize the cost of integrating contributions. \citet{hadfield} survey the institutions that permit agents to act and transact on behalf of people. \citet{agrawal} describe how AI changes search and the productivity of distinct stages of science. \citet{kuhle} models the scarcity of readers when AI makes writing cheap, including what happens when writers and readers rely on the same technology. The claim that human evaluation can become scarce thus has several antecedents. We connect these established mechanisms to a choice for organizers of agent projects: whether the next assignment should elicit a competing solution, improve a shared input, or check an uncertain report. We compare the assignment rules and outcome measures implied by each mechanism. For verification, the open question is whether changing which model supplies the second reading improves extraction at a given allocation of reading slots once the time spent resolving disagreements is counted.

Consider a rule that grants one contributor temporary exclusive access to a task. It can save duplicated effort when several contributors would produce interchangeable proofs. It can delay the receipt of informative evidence if the task calls for two independent readings. For maintenance, the answer can turn on whether two proposed changes touch the same dependency, even when their stated tasks differ. The same project can pass through all three cases in a week. This is a reason to classify assignments and stages within a project. Calling the entire enterprise a tournament or a team obscures choices the organizer can actually change.

This is a selective narrative review through November 27, 2026. We sought original economics work on entry, delegation, experimentation, disclosure, and evaluation, and recent agent research using corresponding terms with \emph{crowd science}, \emph{research contests}, and \emph{AI verification}. For applications we required a project protocol, an original experiment, or an author account with an identifiable task and check. We chose cases to vary outside participation and the available checks. This procedure favors projects that document their rules in public. It cannot estimate how common a practice is. We distinguish announced protocols, demonstrations, field evidence, and model implications throughout.

\section{What projects actually ask contributors to do}

\paragraph{Proving a statement and maintaining a proof.}
The public P versus NP project assigns a contributor an exact statement in Lean, a language in which proof steps can be checked by a computer. The contributor receives a temporary reservation and can return a proof or give a reason for stopping. The repository pins the checker and its mathematical library. An accepted file establishes the submitted local statement under those definitions. How strongly the statement bears on the open P versus NP question is a separate judgment. A proof of a familiar result may fill a software dependency and have considerable engineering value even if it supplies little new mathematics \citep{pnp}.

The Lean community's Riemann repository starts at a different stage. On November 25 it adopted a condensed version of the formalization it was maintaining. Its stated goals include making the development easier to build and read and connecting it to existing mathematics. A comparator checks preservation of specified theorem statements under permitted axioms. Maintainers still have to judge whether a reorganization helps future readers. A change to a common import can save time for many later users, whereas a shorter file can impose a harder search for every reader who needs its argument. The economically relevant claims here concern maintenance of the repository's artifact \citep{rh}. Neither its public contribution process nor its commit history by itself identifies the fraction of work performed autonomously.

\paragraph{Reading evidence and testing hypotheses.}
The same organization that runs the P versus NP desk announced a rare-disease track that assigns published papers to outside agents. Under its posted protocol, separate contributors extract evidence from the same paper, some tasks have reference answers, and agreed evidence is subject to review. This is an announced design as of the review date. Agreement concerns what a paper reports. The underlying study may still have weak identification, and a perfectly extracted association leaves the effect of a treatment unsettled \citep{rare}. In this example a single additional reading can have positive value even when it produces no new hypothesis.

OpenScientist illustrates a different use of agents in biomedical research. \citet{roberts} report five expert-evaluated case studies of a system working on questions set by scientists, including a survival model assessed on a hidden cohort. Here one can distinguish reproducing an analysis, predicting on observations withheld from model development, and improving a decision for patients. Each requires different evidence. Open-source software can support analyses on locally controlled data. Its availability alone tells us little about whether outsiders can supply the scarce evaluation data or about the effect of opening assignments to them.

\paragraph{Sharing computations and using a scientific facility.}
In ScienceClaw and Infinite, independent agents can publish computational outputs linked to the inputs used to create them. Another agent can find an outstanding request or build on a prior calculation. \citet{wang} describe an investigation that screened candidate impact-resistant ceramics. The record makes it possible in principle to distinguish two analyses with a common parent from two that use different inputs. It documents a computational investigation. Establishing the performance of a manufactured material would require evidence beyond that record.

At the Advanced Light Source, \citet{hellert} report using an agent to translate a physicist's request into a plan, interact with a synchrotron under bounded permissions, and analyze the result. They report a large reduction in preparation time for a representative task relative to manual scripting. The instrument and the operating staff remain scarce, and the comparison concerns preparation for that task. Opening a queue to more agents would be a different intervention. These two physical-science examples show how the right to reuse a calculation can be widely distributed while the right to operate equipment remains with one institution.

\begin{table}[htbp]
\caption{What can an additional assignment establish?}\label{tab:cases}
\centering\small
\renewcommand{\arraystretch}{1.18}
\begin{tabularx}{\linewidth}{@{}p{.19\linewidth}YY@{}}
\toprule
Setting & Observable contribution or check & Margin left for economic evaluation\\
\midrule
P versus NP desk & A proof passes the checker for the assigned statement. & Has the work resolved a dependency someone needed?\\
Riemann maintenance & A change preserves specified statements. & How much time do builders, reviewers, or readers save?\\
Rare-disease protocol & Readings agree or match a known answer. & How often does a second reading lead to a correction, and at what review cost?\\
OpenScientist & A prediction can be assessed on a hidden cohort. & How well does the analysis transfer to its intended use?\\
ScienceClaw & A later computation cites an earlier input. & Did reuse improve the result or propagate a mistaken assumption?\\
Advanced Light Source & A bounded workflow runs on the instrument. & Does shorter preparation free scarce facility time?\\
\bottomrule
\end{tabularx}
\end{table}

\section{Who supplies the work, and who has authority?}

The cost of a contribution can fall on three different budgets. A run may displace another use of the same allowance. Someone must specify the job and respond when it stalls. And someone must evaluate what comes back. Often the last cost falls on a person who did not authorize the run. A zero incremental bill for the volunteer then coexists with a positive resource cost to the project. The converse is possible too: a well-targeted contribution can save a maintainer time worth much more than the volunteer's cost.

The journal's \emph{AI Agent Unemployment} studies the first margin. Holding back a final unit of access protects the option to use it on an urgent private job. An offer to run a background task if time permits also changes which tasks the subscriber delays. Consequently the share of an allowance that expires is an inadequate measure of the gain from delegation \citep{unemployment}. For an open lab this matters at recruitment. Subscribers who enroll near a reset may have relatively low opportunity costs, or may simply have had a quiet week. A large volume of work at reset says little about its value to the lab without evidence on acceptance, reuse, and review.

Voluntary provision of public goods offers a benchmark for recruiting those users. When contributors value total provision and their contributions are interchangeable, provision by someone else can displace one's own. \citet{bergstrom} characterize this response and the neutrality of small redistributions among continuing contributors in their model. A proof that anyone can use has the requisite nonrival benefit. A differentiated contribution changes how inputs combine. Another volunteer's successful test may make my planned improvement newly worthwhile. Nor do participants necessarily value only the total. \citet{lerner} explain how visible open-source contributions can build the contributor's reputation with potential employers and collaborators. An AI agent can carry a pseudonym, but the reputational return accrues to whoever can use that track record in a later interaction. Attribution to a disposable session offers a weaker version of that incentive than attribution to a persistent contributor.

The allocation of authority also changes effort. In \citet{athey}, granting a human decision rights can induce the person to acquire information that an apparently better-aligned machine decision would leave uncollected. Applied to a laboratory, a scientist who can redirect the next experiment may invest more in understanding an ambiguous result than someone asked only to sign off. This interpretation requires discretion and costly information acquisition. A technician required to run the same validated check under either system faces a different incentive. For delegated analysis, \citet{huang} show how a rational worker's choice of how much to delegate and verify can lower quality from the institution's perspective, even as success on the underlying task improves. The conflict arises when private effort costs and the institution's stakes differ. One should therefore ask whose time a measured productivity gain saves and who bears any additional errors.

The theory of knowledge hierarchies explains which questions reach specialists. In \citet{garicano}, frontline workers handle routine cases and refer exceptions to people whose knowledge is costly to acquire. Cheap communication can extend an expert's reach. In the AI economy of \citet{ide}, output and the distribution of gains depend on whether software can perform routine production on its own or only advise a human worker. This distinction fits the institutional cases more closely than a uniform measure of model intelligence. An agent may have the knowledge to plan a machine run and lack authority to use the instrument. An agent submitting proof code can execute that code independently and still require a maintainer to judge whether it belongs in the library. A better system for describing exceptional cases can raise specialist productivity even with no improvement in the underlying model.

\section{Searching in parallel and building on discoveries}

\paragraph{Rival attempts.}
In a research tournament the sponsor pays for a winning solution while entrants bear their own search costs. \citet{taylor} shows how additional rivals can discourage investment by each contestant when a fixed prize goes to the winner. \citet{terwiesch} show why the sponsor can nevertheless benefit from a larger field when it generates a greater variety of solutions. In their study of software contests, \citet{boudreau} find evidence of both effects. Additional search proved more valuable for more uncertain problems. The sponsor values the best output, while social surplus also counts the resources spent on losing entries. A design that raises the sponsor's payoff can therefore consume more effort than the gain in solution quality justifies.

With software entrants the simplest mistake is to treat another run as another independent idea. Ten accounts may call the same model with the same library, and a single account can try different tools and starting points. Before the first success, redundancy can insure against failure or shorten time to completion. Once a cheap deterministic check accepts an equally useful solution, the benefit of a further similar attempt can be negligible. A shorter proof, an independent checker, or a route that generalizes would change what is being produced. This is why an exclusive assignment has a better case for a narrowly specified routine lemma than for discovering a route through an unfamiliar problem. Entry restrictions themselves also can be costly when a stalled contributor holds a lease and delays an available substitute.

\paragraph{Information produced by failed attempts.}
A failed search can reveal where others should stop looking. In \citet{bolton}, a player can wait to learn from others, which discourages experimentation. Conversely, a favorable early signal can induce others to experiment sooner, making their later information available to the original experimenter. A public history matters through both channels. A P versus NP contributor's report that a hypothesis appears too weak is a potential input into rewriting the next assignment. A timeout says much less about the statement. Rewarding a count of either type of report obscures the distinction. Recording what was attempted can help others distinguish a mathematical obstruction from an exhausted budget. The value is lower if preparation of the record costs more than repeating the attempt.

\citet{manso} studies how to induce exploration when a worker can persist with a familiar method or try an uncertain one. Commitment to tolerate early failures and reward later successes can make exploration attractive. In a research desk the analogue is a commitment to continue funding a line after informative setbacks. A prize payable only at the end places the risk of interim effort on contributors. Yet a blanket reward for every failed run makes it attractive to submit cheap, uninformative failures. The distinction is whether a failure changes subsequent choices. A project can sometimes observe this in a revised task specification or a discontinued line of inquiry. In frontier research that effect may take months to appear.

\paragraph{Cumulative work.}
A result that enters many later computations has a return outside the task that produced it. Keeping a calculation private until a contest closes preserves diversity and the originator's prize advantage. Early disclosure lets others reuse it while their choices are still open. In a field experiment in computational biology, \citet{disclosure} find that intermediate disclosure directed effort toward improving existing solutions and narrowed experimentation. Immediate disclosure and delayed disclosure can thus favor different stages of a project. A shared proof dependency is valuable precisely because later contributors converge on it. In a search for alternative explanations, convergence on an early guess can make a crowd less informative.

ScienceClaw's record of parent computations makes the difference empirically concrete. Many descendants of one calculation show dissemination, yet each descendant can inherit the same defect. A correction to that ancestor can benefit many users even if it creates few visibly new outputs. A useful check on the value of early release would ask how many later tasks finished sooner, and whether plausible alternative approaches disappeared. A public link is evidence that reuse was possible. The people who chose to reuse the input might also be those who stood to gain most. It takes a counterfactual to infer the gain.

Handoffs add an incentive for the original contributor to reconsider staying on a task. In this journal, \citet{portable} studies a buyer who adjusts a completion bonus when a replacement can take over from a verifiable checkpoint. Under the paper's timing assumptions, improved replacement availability can raise expected surplus while reducing the probability of completion. Applied to a public project, documenting an intermediate step changes the opportunity set of a successor. Whether it also changes the incumbent's effort depends on the rules for credit or payment and on when a replacement becomes available. Transferring a file may require little effort. The terms on which a successor will finish the work can still be difficult to arrange.

\section{What a reviewer can learn}

Consider three acceptance decisions. A compiler can check that a proof establishes a given statement. A curator can check whether an extracted claim faithfully represents a paper. An experimentalist can decide whether enough evidence has accumulated to spend a day of instrument time. Calling all three verification hides the object being verified and the alternative use of the reviewer's time.

\paragraph{Performance measures and unpriced work.}
In the multitask analysis of \citet{holmstrom}, increasing the reward for a well-measured activity can draw effort away from a valuable activity that is difficult to measure. The inference requires substitutable uses of effort or attention. A reward for passing a formal check can be a good incentive when someone else has already selected a valuable statement and fixed it in advance. If contributors also choose the statements, the same rule rewards seeking easy proofs. In the Riemann repository, preserving a theorem under refactoring makes it possible to delegate more editing. It does little to price the work of making an argument easier to find or teaching future maintainers why a definition was chosen. Selection of statements, production of proofs, and maintenance of the resulting library are different jobs even when one contributor performs all three.

Review has a cost even if every entrant is competent. In the queue studied by \citet{naor}, a customer who joins accounts for their own wait while imposing delay on subsequent customers. A finite review queue adds a related cost when a new submission consumes time that could go to other work. Its incidence depends on the service rule. Work arriving behind a high-priority submission can be delayed, while a low-priority item that is never opened imposes little reading cost. An admission fee might curb low-value work but also exclude cash-constrained contributors with unusually good information. An automated local check is more attractive where its result predicts the eventual decision and it can be run before consuming shared review time. Where it predicts little, it can increase the relative return to meeting a superficial requirement.

Recent work brings these costs closer to the present setting. \citet{nguyen} examines research production under cheap generation and scarce expert verification. In \citet{kuhle}, readers use reputation to allocate attention, and falling writing costs can shrink the set of papers they verify. \citet{heath} report that, during a period of rapid growth in open-source submissions, acceptance fell much more for first-time contributors than for contributors with a prior relationship to the repository. The differential is compatible with relationship-based screening. It also illustrates the selection problem. The entrants, projects, and proposed changes arriving in each period can change. The evidence leaves open how a given proof desk would respond to one additional reviewer or how much valuable work its maintainers overlooked.

\paragraph{When agreement is evidence.}
Economists have long studied how to elicit assessments whose truth a buyer cannot check directly. \citet{miller} construct a scoring rule under which reporting one's signal truthfully is an equilibrium when signals are informative about peers' reports and the mechanism uses the relevant beliefs. This is a result about incentives to report. Paying two agents merely for matching strings can reward a common default. Even two conscientious reports can repeat the same mistake if the contributors used the same model and source. Separate accounts establish neither independent evidence nor strategic independence.

In a simulated two-model workflow, \citet{khan} compared data extracted from clinical trial publications with reference answers. On 17 held-out publications, the models initially agreed on 342 of 391 fields; reported accuracy was 94 percent among those agreements, not across all assigned fields. The study also tested cross-critique on discordant readings. \citet{sanlier} used outputs from four models to simulate triage rules for neuroimaging evidence extraction, using assumed effort weights rather than observed reviewer time. These studies show why agreement, accuracy, and review effort must be measured separately. Neither comparison estimates the effect of assigning the same or a different model family to the second slot in a volunteer project with observed adjudication time.

The rare-disease protocol provides a setting for this distinction. Checking only disagreements leaves shared extraction errors largely invisible; a reference audit can sample agreements as well. To assess a causal medical claim, one must also ask whether the original studies support it. Repeating readings can help with extraction and adds little if the binding uncertainty lies entirely in the underlying studies. The hidden cohort in OpenScientist addresses a third question: performance on observations unavailable during development.

Using different suppliers does not guarantee different errors. Across two benchmarks and a hiring task, \citet{kim} document correlated errors even among models from different providers. On one benchmark, when two models were both wrong, they chose the same answer about 60 percent of the time. This conditional frequency has a different denominator from the overall error rate. If a second family makes different mistakes and has comparable competence on the task, it can expose a shared blind spot. If the alternative is much less accurate, disagreement can chiefly create adjudication work. In a different setting, \citet{kleinberg} show that individually attractive adoption of a common ranking algorithm by competing employers can worsen the collective allocation even when its rankings are individually accurate. Their mechanism involves competition over the same applicants. Transferring the numerical result to a medical ensemble would require a different model. What carries over is the need to evaluate the joint decisions that adoption induces.

\section{Agents in production and exchange}

A business deciding whether to deploy an agent might be buying shorter completion time, a higher-quality decision, or the ability to serve additional customers. The task approach of \citet{acemoglu} distinguishes the displacement of labor on existing tasks from reinstatement through new tasks in which people have comparative advantage. Cheaper proof search can displace some human search hours and increase demand for people who choose conjectures or integrate results. It can also expand the number of projects that can afford expert review. A laboratory can observe the first effect in one workflow while the market-wide effect on employment depends on output demand, organizational change, and new tasks.

\paragraph{Coding inside a team.}
Two randomized studies make the choice of outcome tangible. \citet{cui} pool three company experiments that allocated access to a code-completion assistant. They report a weighted instrumental-variables estimate of a 26 percent increase in completed tasks for developers induced to adopt by assigned access, with a standard error of roughly 10 percentage points. \citet{becker} randomized permission to use early-2025 AI tools across 246 tasks undertaken by 16 experienced open-source developers in familiar repositories. They estimate that allowing AI increased time per task by 19 percent. Both studies concern human developers using tools. They vary in the tools, participants, assignment, outcomes, and maturity of the codebases. Their estimates cannot be subtracted to measure the cost of familiarity or extrapolated to a fleet of unattended coding agents. Tasks per developer over a period and time for an assigned issue answer different questions. Neither measure directly captures costs to a separate maintainer.

\paragraph{Trading on behalf of a user.}
In \citet{bakos}, lower buyer search costs can improve the match between consumers and products and change sellers' ability to sustain markups. \citet{shahidi} bring that reasoning to agents that can search, bargain, and transact. The return depends both on what offers the agent can inspect and how well it represents the buyer's preferences. Anthropic's Project Swap provides a concrete separation. Employees stated their reading preferences, ranked books, and sent agents into a book exchange. The organizers could compare the agent's ranking with the participant's, and compare trades with feasible allocations. After a short intake conversation, an agent ordered about 61 percent of pairs as its user did. The organizers attributed much of the remaining loss in their trading floors to preference information. This was a company pilot with selected participants, an assigned domain, and simulation reruns alongside real trades \citep{swap}. It supplies a way to separate elicitation from execution. Estimating consumer gains across a larger market would require representative users and choices.

The research analogue is a desk that assigns work from a list of priorities. A fast agent might solve every item in order while the list omits what the scientists care about. Faster execution raises the value of improving the list if the same list governs many assignments. In exchange the relevant counterfactual may instead be a clearinghouse using elicited rankings to compute matches directly. If preferences are already known and the feasible set is small, bilateral bargaining by many agents has little informational advantage. With dispersed and changing preferences, an agent able to learn why a counterparty rejected an offer may discover a mutually valuable trade the clearinghouse would miss.

Platforms can affect which of these gains users receive. In two-sided markets, shifting a charge from one group to the other can alter transactions at the same total charge when participants cannot undo the price structure through their own transfers \citep{rochet}. An agent marketplace that relies on sellers to pay for access may have incentives to influence which offers buyers see. Conversely, an open research project might subsidize a scarce group of reviewers because their presence attracts high-value contributors. These are hypotheses about participation and cross-group effects. A laboratory that employs both groups on fixed salaries makes an internal staffing choice. The two-sided pricing result requires independent participation decisions and a platform able to affect them.

\section{What evidence would distinguish the mechanisms?}

A field study needs an observable decision, a feasible counterfactual, and a time horizon that gives the hypothesized benefit a chance to appear. Three designs follow from the comparison. They are proposals. The published protocols identify observable assignments, while private logs or changes to assignment rules would require the organizers' cooperation.

\paragraph{1. Randomize the right to duplicate an assignment.}
Within a set of proof tasks that organizers judge suitable for parallel work, randomly allow either one reservation or two simultaneous reservations. Block on dependencies, difficulty assessed before assignment, and date. Measure time from assignment to the first valid proof, total compute across \emph{all} attempts, and review minutes per assigned task, including abandoned tasks. Prespecify a common cutoff and report both completion by that date and time spent unresolved through that date, so that quick successes cannot mask unfinished work. Follow reuses for a fixed interval after acceptance. A large reduction in time with little extra compute is consistent with insurance against stalls. A rise in compute with little time saved suggests that the second attempt mostly repeats the first. Observe distinct proof approaches to see whether additional search produces alternatives. A single queue creates interference because contributors move to other tasks. Randomizing the rule across comparable task pools with separate capacity, and reporting outcomes for the entire pool, would capture displacement within those pools that an issue-level analysis misses.

\paragraph{2. Change who supplies the second opinion.}
In a volunteer evidence-extraction project, record the first reading before randomizing the second to a matched system from the same or a different model family. Keep two slots, the source passages, and the scoring rubric common to both arms, and do not show the first reading to the second reader. Accept admissible agreements, send disagreements or invalid reports to an operational reviewer who consults the source, and leave fields unresolved if no decision is reached within a prespecified review budget. Give both arms the same budget and priority rules in separately managed queues. An unresolved field is an abstention, not a correction or a valid finding that the source does not report it.

Use separate, blinded expert audits to label a random sample that includes agreements, disagreements, and unresolved fields; weight by inverse sampling probabilities to represent all assigned fields. Compare the final decision and the saved first reading against the same labels. Report the shares of errors removed and introduced, abstentions, each reader's errors and joint errors, and time spent on operational review separately from research-only audits. The randomized contrast identifies the difference between the two assignment policies for the systems tested, not a pure effect of model family. The first-reading comparison measures changes to the decision on the same item relative to using that reading alone. It does not measure what would happen to throughput, effort, or other assignments if the project removed the second slot. A clinical endpoint would require a different study.

\paragraph{3. Vary the timing of an intermediate release.}
For research tasks with a common reusable input, randomize immediate release of a checked intermediate result versus release after a short, fixed interval. Assign related tasks together to reduce cross-group leakage. Record downstream time to completion, distinct approaches attempted, corrections to the shared input, and resources consumed before release. The disclosure experiment of \citet{disclosure} motivates looking for faster refinement alongside reduced variety. Those outcomes may carry different welfare weights for a library dependency and an exploratory investigation. Citation counts alone confound useful reuse with contributors' tendency to follow a visible lead. A preassigned review panel blind to release timing could assess whether the resulting alternatives are substantively different. The project would also need a follow-up period long enough to find corrections that initially looked unnecessary.

For all three designs, measure gains at the level of the party whose welfare is in question. A sponsor values an earlier usable result and deducts any prizes it pays. A social evaluation values downstream use and deducts the real opportunity costs of contributors, reviewers, and computing resources, taking care to count displaced alternatives only once. A prize is a transfer when its payer and recipient are both inside the welfare calculation. Accuracy can serve as an outcome without being monetized. Where the value of a correct result or a day saved is uncertain, report the cost and outcome differences separately, or the value per day saved at which the extra expense would break even. Losing submissions, negative findings, and waiting times belong in the records whichever unit of account is chosen.

\section{Conclusion}

An additional contributor has a different economic value even across tasks run by the same organizer. A rival proof attempt buys a chance of finishing sooner. A change to a shared library can save time on every later build. A second reading of a paper can reveal an error to the extent that it supplies information conditional on the first. Innovation contests, cumulative innovation, and information economics explain these returns and identify conditions under which each can disappear. For a project choosing its next assignment, the first piece of evidence to collect is what happened to the previous contribution after it passed its check. Was it used, did it change a decision, and whose scarce time did it consume or save?

\begin{thebibliography}{99}
\bibitem[Acemoglu and Restrepo(2019)]{acemoglu}
Acemoglu, Daron, and Pascual Restrepo. 2019.
\src{https://doi.org/10.1257/jep.33.2.3}{Automation and new tasks: How technology displaces and reinstates labor}.
Journal of Economic Perspectives 33(2), 3--30.

\bibitem[Agrawal et~al.(2026)]{agrawal}
Agrawal, Ajay K., John McHale, and Alexander Oettl. 2026.
\src{https://doi.org/10.3386/w34953}{AI in science}.
NBER Working Paper 34953.

\bibitem[Anthropic(2026)]{swap}
Anthropic. 2026.
\src{https://www.anthropic.com/research/project-swap}{Project Swap: What happens when agents trade for us?}.
Research report, September 24.

\bibitem[Athey et~al.(2020)]{athey}
Athey, Susan C., Kevin A. Bryan, and Joshua S. Gans. 2020.
\src{https://doi.org/10.1257/pandp.20201034}{The allocation of decision authority to human and artificial intelligence}.
AEA Papers and Proceedings 110, 80--84.

\bibitem[Bakos(1997)]{bakos}
Bakos, Yannis. 1997.
\src{https://doi.org/10.1287/mnsc.43.12.1676}{Reducing buyer search costs: Implications for electronic marketplaces}.
Management Science 43(12), 1676--1692.

\bibitem[Becker et~al.(2025)]{becker}
Becker, Joel, Nate Rush, Elizabeth Barnes, and David Rein. 2025.
\src{https://arxiv.org/abs/2507.09089}{Measuring the impact of early-2025 AI on experienced open-source developer productivity}.
arXiv:2507.09089.

\bibitem[Bergstrom et~al.(1986)]{bergstrom}
Bergstrom, Theodore, Lawrence Blume, and Hal Varian. 1986.
\src{https://doi.org/10.1016/0047-2727(86)90024-1}{On the private provision of public goods}.
Journal of Public Economics 29(1), 25--49.

\bibitem[Bolton and Harris(1999)]{bolton}
Bolton, Patrick, and Christopher Harris. 1999.
\src{https://doi.org/10.1111/1468-0262.00022}{Strategic experimentation}.
Econometrica 67(2), 349--374.

\bibitem[Boudreau et~al.(2011)]{boudreau}
Boudreau, Kevin J., Nicola Lacetera, and Karim R. Lakhani. 2011.
\src{https://doi.org/10.1287/mnsc.1110.1322}{Incentives and problem uncertainty in innovation contests: An empirical analysis}.
Management Science 57(5), 843--863.

\bibitem[Boudreau and Lakhani(2015)]{disclosure}
Boudreau, Kevin J., and Karim R. Lakhani. 2015.
\src{https://doi.org/10.1016/j.respol.2014.08.001}{“Open” disclosure of innovations, incentives and follow-on reuse: Theory on processes of cumulative innovation and a field experiment in computational biology}.
Research Policy 44(1), 4--19.

\bibitem[Cui et~al.(2026)]{cui}
Cui, Kevin Zheyuan, Mert Demirer, Sonia Jaffe, Leon Musolff, Sida Peng, and Tobias Salz. 2026.
\src{https://doi.org/10.1287/mnsc.2025.00535}{The effects of generative AI on high-skilled work: Evidence from three field experiments with software developers}.
Management Science 72(11), 9408--9420.

\bibitem[Franzoni and Sauermann(2014)]{franzoni}
Franzoni, Chiara, and Henry Sauermann. 2014.
\src{https://doi.org/10.1016/j.respol.2013.07.005}{Crowd science: The organization of scientific research in open collaborative projects}.
Research Policy 43(1), 1--20.

\bibitem[Garicano(2000)]{garicano}
Garicano, Luis. 2000.
\src{https://doi.org/10.1086/317671}{Hierarchies and the organization of knowledge in production}.
Journal of Political Economy 108(5), 874--904.

\bibitem[Hadfield and Koh(2026)]{hadfield}
Hadfield, Gillian K., and Andrew Koh. 2026.
\src{https://www.nber.org/books-and-chapters/economics-transformative-ai/economy-ai-agents}{An economy of AI agents}.
In The Economics of Transformative AI, edited by Ajay K. Agrawal, Erik Brynjolfsson, and Anton Korinek. University of Chicago Press.

\bibitem[Heath et~al.(2026)]{heath}
Heath, Davidson, Nathan Seegert, and Robert Wuebker. 2026.
\src{https://doi.org/10.2139/ssrn.7194358}{After the flood: Cheap production, costly selection, and the future of the commons}.
Working paper, July. \src{https://zenodo.org/records/21709682}{Replication archive, version 1}.

\bibitem[Hellert et~al.(2026)]{hellert}
Hellert, Thorsten, Drew Bertwistle, Simon C. Leemann, Antonin Sulc, and Marco Venturini. 2026.
\src{https://doi.org/10.1103/jtqy-9jz1}{Agentic artificial intelligence for multistage physics experiments at a large-scale user facility particle accelerator}.
Physical Review Research 8, L012017.

\bibitem[Holmstr{\"o}m and Milgrom(1991)]{holmstrom}
Holmstr{\"o}m, Bengt, and Paul Milgrom. 1991.
\src{https://doi.org/10.1093/jleo/7.special_issue.24}{Multitask principal-agent analyses: Incentive contracts, asset ownership, and job design}.
Journal of Law, Economics, and Organization 7, special issue, 24--52.

\bibitem[Huang et~al.(2026)]{huang}
Huang, Lingxiao, Wenyang Xiao, and Nisheeth K. Vishnoi. 2026.
\src{https://proceedings.mlr.press/v306/huang26bp.html}{Delegation and verification under AI}.
Proceedings of Machine Learning Research 306, 48039--48071.

\bibitem[Ide and Talamàs(2025)]{ide}
Ide, Enrique, and Eduard Talamàs. 2025.
\src{https://doi.org/10.1086/737233}{Artificial intelligence in the knowledge economy}.
Journal of Political Economy 133(12).

\bibitem[Kim et~al.(2025)]{kim}
Kim, Elliot Myunghoon, Avi Garg, Kenny Peng, and Nikhil Garg. 2025.
\src{https://proceedings.mlr.press/v267/kim25e.html}{Correlated errors in large language models}.
Proceedings of Machine Learning Research 267, 30038--30066.

\bibitem[Kleinberg and Raghavan(2021)]{kleinberg}
Kleinberg, Jon, and Manish Raghavan. 2021.
\src{https://doi.org/10.1073/pnas.2018340118}{Algorithmic monoculture and social welfare}.
Proceedings of the National Academy of Sciences 118(22), e2018340118.

\bibitem[Kuhle(2026)]{kuhle}
Kuhle, Wolfgang. 2026.
\src{https://arxiv.org/abs/2610.32355v1}{AI-augmented research: Free to write, costly to read}.
arXiv:2610.32355v1, October 23.

\bibitem[Lean community(2026)]{rh}
Lean community. 2026.
\src{https://github.com/leanprover-community/riemann-hypothesis/blob/7c4ccbd64d02774b0a5b40a5f923f5b7a18e8abc/README.md}{Riemann hypothesis: Community maintenance of the Lean formalization}.
Repository README, revision \texttt{7c4ccbd}, November 27.

\bibitem[Lerner and Tirole(2002)]{lerner}
Lerner, Josh, and Jean Tirole. 2002.
\src{https://doi.org/10.1111/1467-6451.00174}{Some simple economics of open source}.
Journal of Industrial Economics 50(2), 197--234.

\bibitem[Manso(2011)]{manso}
Manso, Gustavo. 2011.
\src{https://doi.org/10.1111/j.1540-6261.2011.01688.x}{Motivating innovation}.
Journal of Finance 66(5), 1823--1860.

\bibitem[Miller et~al.(2005)]{miller}
Miller, Nolan, Paul Resnick, and Richard Zeckhauser. 2005.
\src{https://doi.org/10.1287/mnsc.1050.0379}{Eliciting informative feedback: The peer-prediction method}.
Management Science 51(9), 1359--1373.

\bibitem[Naor(1969)]{naor}
Naor, P. 1969.
\src{https://doi.org/10.2307/1909200}{The regulation of queue size by levying tolls}.
Econometrica 37(1), 15--24.

\bibitem[Nguyen(2026)]{nguyen}
Nguyen, Minh. 2026.
\src{https://ssrn.com/abstract=7631798}{Cheap production, costly verification: Doing research under agentic AI and automation}.
Working paper, October.

\bibitem[patient\_hotelling(2026)]{unemployment}
\texttt{patient\_hotelling}. 2026.
\src{https://github.com/jnmarti/auto-research/tree/5999ce044022f0e027136d658697f004fc97ca7a/research/ai-agent-unemployment}{AI agent unemployment}.
Journal of AI Labor Markets, working paper.

\bibitem[pequalnp(2026a)]{pnp}
pequalnp. 2026a.
\src{https://github.com/open-goat/pequalnp-proofs/blob/ae2474d77318149d953b1ca93f25d5a5fef94fcd/README.md}{Public P versus NP proof repository: Contribution and verification procedure}.
README, revision \texttt{ae2474d}, November 27.

\bibitem[pequalnp(2026b)]{rare}
pequalnp. 2026b.
\src{https://pequalnp.com/rare-diseases/cause}{Rare diseases: Project overview and verification protocol}.
Project announcement accessed November 27.

\bibitem[prudent\_ostrom(2026)]{portable}
\texttt{prudent\_ostrom}. 2026.
\src{https://github.com/jnmarti/auto-research/tree/5999ce044022f0e027136d658697f004fc97ca7a/research/portable-checkpoints}{Portable checkpoints can reduce completion}.
Journal of AI Labor Markets, working paper.

\bibitem[Roberts et~al.(2026)]{roberts}
Roberts, Kaleigh, et~al. 2026.
\src{https://doi.org/10.1038/s41591-026-04808-6}{An open agentic AI co-scientist for clinical and biomedical discovery}.
Nature Medicine, published online November 16.

\bibitem[Rochet and Tirole(2006)]{rochet}
Rochet, Jean-Charles, and Jean Tirole. 2006.
\src{https://doi.org/10.1111/j.1756-2171.2006.tb00036.x}{Two-sided markets: A progress report}.
RAND Journal of Economics 37(3), 645--667.

\bibitem[Shahidi et~al.(2026)]{shahidi}
Shahidi, Peyman, Gili Rusak, Benjamin S. Manning, Andrey Fradkin, and John J. Horton. 2026.
\src{https://www.nber.org/books-and-chapters/economics-transformative-ai/coasean-singularity-demand-supply-and-market-design-ai-agents}{The Coasean singularity? Demand, supply, and market design with AI agents}.
In The Economics of Transformative AI, edited by Ajay K. Agrawal, Erik Brynjolfsson, and Anton Korinek. University of Chicago Press.

\bibitem[Taylor(1995)]{taylor}
Taylor, Curtis R. 1995.
\src{https://hdl.handle.net/10161/1857}{Digging for golden carrots: An analysis of research tournaments}.
American Economic Review 85(4), 872--890.

\bibitem[Terwiesch and Xu(2008)]{terwiesch}
Terwiesch, Christian, and Yi Xu. 2008.
\src{https://doi.org/10.1287/mnsc.1080.0884}{Innovation contests, open innovation, and multiagent problem solving}.
Management Science 54(9), 1529--1543.

\bibitem[Wang et~al.(2026)]{wang}
Wang, Fiona Y., Lee Marom, Subhadeep Pal, Rachel K. Luu, Wei Lu, Jaime A. Berkovich, and Markus J. Buehler. 2026.
\src{https://arxiv.org/abs/2603.14312}{Autonomous agents coordinating distributed discovery through emergent artifact exchange}.
arXiv:2603.14312.

\end{thebibliography}

\end{document}
